\usepackage{xstring}

% macro vérifiées


% #1 = width, #2=height, #3=list de c, d
% c pour /
% d pour \
\newcommand{\slant}[3]{
	\setcounter{col}{0}
	\setcounter{row}{0}
	\begin{scope}[shift={(0,#2)},yscale=-1]
	\foreach \it in #3 {
		\ifthenelse{ \value{col}= #1 }{
			\stepcounter{row};
			\setcounter{col}{0};
		}{}
		\ifthenelse{ \equal{\it}{c} }{
			\draw[line width=4pt] ({\value{col}},{\value{row}+1}) --++ (1,-1);
		}{
		  \ifthenelse{ \equal{\it}{d} }{
			\draw[line width=4pt] ({\value{col}},{\value{row}}) --++ (1,1);
		  }{}
		}
		\stepcounter{col};
	}
	\end{scope}
}

% #1 = width, #2=height, #3=list, #4=size du texte, #4=rayon des îles
\newcommand{\islands}[5]{
    \setcounter{col}{0}
	\setcounter{row}{0}
	\begin{scope}[shift={(0,#2)},yscale=-1]
	\foreach \it in #3 {
		\ifthenelse{ \value{col}= #1 }{
			\stepcounter{row};
			\setcounter{col}{0};
		}{}
		\ifthenelse{ \equal{\it}{}}{}{
		    \draw[fill=white] ({\value{col}},{\value{row}+1}) circle(#5) node[scale=#4]{\it};
		}
		\stepcounter{col};
	}
	\end{scope}
}


% #1 = width, #2=height, #3=list, #4=rayon
\newcommand{\showCircleList}[4]{
	\setcounter{col}{0}
	\setcounter{row}{0}
	\begin{scope}[shift={(.5,{#2-.5})},yscale=-1]
	\foreach \it in #3 {
		\ifthenelse{ \value{col}= #1 }{
			\stepcounter{row};
			\setcounter{col}{0};
		}{}
		\ifthenelse{ \equal{\it}{B} }{
            % disque blanc    		
    		\draw[fill=white] (\value{col},\value{row}) circle(#4);%
		}{
			\ifthenelse{ \equal{\it}{N} }{
				% disque noir
				\draw[fill=black] (\value{col},\value{row}) circle(#4);%
			}{
				% rien
			}
		}
		\stepcounter{col};
	}
	\end{scope}
}






% à vérifier

\newcommand\sudokuGrid{
\draw[line width=3pt] (0,0) rectangle (9,9);
\draw[line width=2pt] (0,0) grid[step=3] (9,9);
\draw[line width=1pt] (0,0) grid[step=1] (9,9);
}

\newcommand\sudokuGridP{
\draw[line width=3pt] (0,0) rectangle (4,4);
\draw[line width=2pt] (0,0) grid[step=2] (4,4);
\draw[line width=1pt] (0,0) grid[step=1] (4,4);
}

\newcommand\plainGrid[1]{
\draw[line width=3pt] (0,0) rectangle (#1,#1);
\draw[line width=1pt] (0,0) grid[step=1] (#1,#1);
}

\newcommand\lightPlainGrid[1]{
\draw[line width=2pt] (0,0) rectangle (#1,#1);
\draw[line width=0.2pt] (0,0) grid[step=1] (#1,#1);
}


\newcommand\dottedGrid[1]{
\draw[dotted] (0,0) rectangle (#1,#1);
\draw[dotted] (0,0) grid[step=1] (#1,#1);
}

\newcommand\dottedGridRect[2]{
\draw[dotted] (0,0) rectangle (#1,#2);
\draw[dotted] (0,0) grid[step=1] (#1,#2);
}



\newcommand{\quadruple}[1]{
\foreach \x/\y\lab in #1 {
	\draw[fill=white] (\x,\y) circle (.35);
	\draw (\x,\y) node[scale=.7]{\lab};
}
}


\newcommand{\xyBlackSquares}[2]{
\begin{scope}[shift={(0,#1)}, yscale=-1]
\foreach \x/\y in #2 {
	\fill[black] (\x,\y) rectangle ++ (1,1);
}
\end{scope}
}

\newcommand{\xyCircle}[3]{
\foreach \x/\y\/z in #1 {
	\draw[black] (\x,\y) cirle(#2) node[scale=#3] {\z};
}
}



\newcommand{\bridgesCor}[2]{
	\begin{scope}[shift={(0,#1)},yscale=-1]	
	\foreach \lineA/\colA/\lineB/\colB/\val in #2 {
		\ifthenelse{\val = 1}{
			\draw[thick](\colA,\lineA) -- (\colB,\lineB);
		}{
			\ifthenelse{\val = 2}{
			    \draw[thick, double, double distance=1mm](\colA,\lineA) -- (\colB,\lineB);		
			}{}
		}
	}
	\end{scope}
}







% r pour |
% s pour -
% t pour gauche - haut
% u pour gauche - bas
% v pour haut - droite
% w pour bas - droite
\newcommand{\showBlackList}[3]{
	\setcounter{col}{0}
	\setcounter{row}{0}
	\begin{scope}[shift={(0,#1)},yscale=-1]
	\foreach \it in #2 {
		\ifthenelse{ \value{col}= #1 }{
			\stepcounter{row};
			\setcounter{col}{0};
		}{}
		\ifthenelse{ \equal{\it}{b} }{
			\fill[#3] ({\value{col}},{\value{row}}) rectangle ++ (1,1);
		}{
		  \ifthenelse{ \equal{\it}{r} }{
			\fill[#3] ({\value{col}+.1},{\value{row}}) rectangle ++ (.8,1);
		  }{
		    \ifthenelse{ \equal{\it}{s} }{
			  \fill[#3] ({\value{col}},{\value{row}+.1}) rectangle ++ (1,.8);
		    }{
		    \ifthenelse{ \equal{\it}{t} }{
			  \fill[#3] ({\value{col}+.1},{\value{row}}) -|++ (.8,.9) -|++ (-.9,-.8) -| cycle;
		    }{
		    \ifthenelse{ \equal{\it}{u} }{
			  \fill[#3] ({\value{col}},{\value{row}+.1}) -|++ (.9,.9) -|++ (-.8, -.1) -| cycle;
		    }{
		    \ifthenelse{ \equal{\it}{v} }{
			  \fill[#3] ({\value{col}+.1},{\value{row}}) -|++ (.8,.1) -|++ (.1,.8) -| cycle;
		    }{
		    \ifthenelse{ \equal{\it}{w} }{
			  \fill[#3] ({\value{col}+.1},{\value{row}+.1}) -|++ (.9,.8) -|++ (-.1,.1) -| cycle;
		    }{
		    }
		    }
		    }
		    }
		    }
		  }
		}

		\stepcounter{col};
	}
	\end{scope}
}

\newcommand{\showListRaw}[3]{
	\setcounter{col}{0}
	\setcounter{row}{0}
	\begin{scope}[shift={(.5,.5)}]
	\foreach \it in #2 {
		\ifthenelse{ \value{col}= #1 }{
			\stepcounter{row};
			\setcounter{col}{0};
		}{}
		\ifthenelse{ \equal{\it}{b} }{
			\fill[black!50] ({\value{col}-.5},{\value{row}-.5}) rectangle ++ (1,1);
		}{
			\ifthenelse{ \equal{\it}{+} }{
				\node[anchor=center, scale=#3] at (\value{col},\value{row}) {$+$};
			}{
				\node[anchor=center, scale=#3] at (\value{col},\value{row}) {\it};			
			}
		}

		\stepcounter{col};
	}
	\end{scope}
}

% #1 = width, #2=height, #3 = list, #4=text size
\newcommand{\showList}[4]{
	\setcounter{col}{0}
	\setcounter{row}{0}
	\begin{scope}[shift={(.5,{#2-.5})},yscale=-1]
	\foreach \it in #3 {
		\ifthenelse{ \value{col}= #1 }{
			\stepcounter{row};
			\setcounter{col}{0};
		}{}
		\ifthenelse{ \equal{\it}{b} }{
			\fill[black!50] ({\value{col}-.5},{\value{row}-.5}) rectangle ++ (1,1);
		}{
			\ifthenelse{ \equal{\it}{+} }{
				\node[anchor=center, scale=#4] at (\value{col},\value{row}) {$+$};
			}{
				\node[anchor=center, scale=#4] at (\value{col},\value{row}) {\it};			
			}
		}

		\stepcounter{col};
	}
	\end{scope}
}



% connexion horizontale
\newcommand{\connectListH}[2]{%#size ; #2 list
	\setcounter{col}{0}
	\setcounter{row}{0}
	\begin{scope}[shift={(0,{#1-.5})},yscale=-1]
	\foreach \it in #2 {
		\ifthenelse{ \value{col}= #1 }{
			\stepcounter{row};
			\setcounter{col}{0};
		}{}
		\ifthenelse{ \equal{\it}{3} }{
            % trait complet
    		\draw[line width=1pt] (\value{col},\value{row}) --++ (1,0);%
		}{
			\ifthenelse{ \equal{\it}{2} }{
				% début seulement
				\draw[line width=1pt] (\value{col},\value{row}) --++ (.5,0);%
			}{
                \ifthenelse{ \equal{\it}{1} }{
                    % fin seulement
                    \draw[line width=1pt] (\value{col},\value{row}) ++ (.5,0) --++(.5,0);%
                }{				
				    % rien
				}
			}
		}

		\stepcounter{col};
	}
	\end{scope}
}

% connexion horizontale
\newcommand{\undead}[3]{%#size ; #2 list
	\setcounter{col}{0}
	\setcounter{row}{0}
	\begin{scope}[shift={(0.5,{#1-.5})},yscale=-1]
	\foreach \it in #2 {
		\ifthenelse{ \value{col}= #1 }{
			\stepcounter{row};
			\setcounter{col}{0};
		}{}
		\ifthenelse{ \equal{\it}{L} }{
            % miroir gauche
    		\draw[line width=2pt] (\value{col},\value{row}) ++ (-.3,-.3) --++ (.6,.6);%
		}{
		\ifthenelse{ \equal{\it}{R} }{
    		% miroir droit
			\draw[line width=2pt] (\value{col},\value{row}) ++ (-.3,.3) --++ (.6,-.6);%
		}{
            % symbole
            \draw (\value{col},\value{row}) node[scale=#3]{\it};%
		}
		}

		\stepcounter{col};
	}
	\end{scope}
}

% connexion verticale
\newcommand{\connectListV}[2]{%#size ; #2 list
	\setcounter{col}{0}
	\setcounter{row}{0}
	\begin{scope}[shift={(.5,#1)},yscale=-1]
	\foreach \it in #2 {
		\ifthenelse{ \value{col}= #1 }{
			\stepcounter{row};
			\setcounter{col}{0};
		}{}
		\ifthenelse{ \equal{\it}{3} }{
            % trait complet
    		\draw[line width=1pt] (\value{col},\value{row}) --++ (0,1);%
		}{
			\ifthenelse{ \equal{\it}{2} }{
				% début seulement
				\draw[line width=1pt] (\value{col},\value{row}) --++ (0,.5);%
			}{
                \ifthenelse{ \equal{\it}{1} }{
                    % fin seulement
                    \draw[line width=1pt] (\value{col},\value{row}) ++ (0,.5) --++(0,.5);%
                }{				
				    % rien
				}
			}
		}

		\stepcounter{col};
	}
	\end{scope}
}



\newcommand{\sideItems}[4]{
% Items écrit sur les bords HBGD
\def\w{#1}
\def\h{#2}
\foreach \sideValues [count=\side] in #3 {
	\ifthenelse{\side = 1}{
		\foreach \it [count=\i] in \sideValues {
			\draw ({\i-.5},\h) node[above, scale=#4]{\it};
		}
	}{}
	\ifthenelse{\side = 2}{
		\foreach \it [count=\i] in \sideValues {
			\draw ({\i-.5},0) node[below, scale=#4]{\it};
		}
	}{}
	\ifthenelse{\side = 3}{
		\foreach \it [count=\i] in \sideValues {
			\draw (0,{-\i+.5+\h}) node[left, scale=#4]{\it};
		}
	}{}	
	\ifthenelse{\side = 4}{
		\foreach \it [count=\i] in \sideValues {
			\draw (\w, {-\i+.5+\h}) node[right, scale=#4]{\it};
		}
	}{}	

}
}

\newcommand\tags[2]{
\foreach \x/\y/\t in  #1 {
	\node[anchor=north west, scale=#2, fill=white] at (\x,\y) {\t};
}
}

\newcommand\rectanglesForKiller[1]{
\def\m{.1}; % marge pour les zones intérieures
\foreach \x/\y/\a/\b in  #1 {
	\draw[line width=.5pt] ({\x+\m},{\y-\m}) rectangle ({\x+\a-\m},{\y-\b+\m} );
}
}

\newcommand\drawZone[5]{
	\begin{scope}[shift={(#1,#2)}, rotate=#3]
		\draw[line width=#5 pt] plot coordinates #4 -- cycle;
	\end{scope}
}

\newcommand\dedalopath[2]{
	\begin{scope}[shift={(-.5,-.5)}]	
		\draw[line width=5pt, #1,*-*] plot coordinates #2;
	\end{scope}
}

\newcommand\thermopath[2]{
	\StrBetween{#1}{(}{,}[\xdeb]
	\StrBetween{#1}{,}{)}[\ydeb]
	\begin{scope}[shift={(-.5,-.5)}]
		\draw[line width=10pt, black!30] (\xdeb , \ydeb) -- plot coordinates #2;
		\fill[black!30] (\xdeb , \ydeb) circle (.4);
	\end{scope}
}

\newcommand\thermopathBlack[2]{
	\StrBetween{#1}{(}{,}[\xdeb]
	\StrBetween{#1}{,}{)}[\ydeb]
	\begin{scope}[shift={(-.5,-.5)}]
		\draw[line width=0.5cm, black] (\xdeb , \ydeb) -- plot coordinates #2;
		\fill[black] (\xdeb , \ydeb) circle (.4);
	\end{scope}
}


\newcommand\thermopathequal[2]{
	\StrBetween{#1}{(}{,}[\xdeb]
	\StrBetween{#1}{,}{)}[\ydeb]
	\begin{scope}[shift={(-.5,-.5)}]
		\draw[line width=8pt, black!30, ->, >=stealth] (\xdeb , \ydeb) -- plot coordinates #2;
		\draw[line width=8pt, black!30, fill=white] (\xdeb , \ydeb) circle (.4);
	\end{scope}
}

\newcommand\inequalGrid[1]{
\foreach \x in {1,2,...,#1} {
	\foreach \y in {1,2,...,#1}{
		\draw ({\x-1+.2},{\y-1+.2}) rectangle ++ (.6,.6);
	}
}
}

\newcommand\inequalBS[1]{
% black square
\foreach \x/\y in #1 {
	\fill[black] ({\x+.2},{\y+.2}) rectangle ++ (.6,.6);
}
}

\newcommand\adjacentH[2]{
% Placement des symboles horizontaux
\setcounter{col}{1}
\setcounter{row}{0}
\begin{scope}[shift={(0,{#1-.5})}, yscale=-1]
% 0 pour rien sinon |
\foreach \it in #2 {
	\ifthenelse{ \value{col}=#1 }{
		\stepcounter{row};
		\setcounter{col}{1};
	}{}
	\ifthenelse{ \it=0 }{
	}{
		\draw[line width=3pt] (\value{col},\value{row}+.2) --++ (0,-.4);
	}
	\stepcounter{col};
}
\end{scope}
}

\newcommand\adjacentV[2]{
% Placement des symboles verticaux
\setcounter{col}{0}
\setcounter{row}{0}
\begin{scope}[shift={(.5,{#1-1})}, yscale=-1]
% 0 pour rien sinon |
\foreach \it in #2 {
	\ifthenelse{ \value{col}=#1 }{
		\stepcounter{row};
		\setcounter{col}{0};
	}{}
	\ifthenelse{ \it=0 }{
	}{
		\draw[line width=3pt] ({\value{col}-.2},\value{row}) --++ (.4,0);
	}
	\stepcounter{col};
}
\end{scope}
}


\newcommand\inequalH[2]{
% Placement des symboles horizontaux
\setcounter{col}{1}
\setcounter{row}{0}
\begin{scope}[shift={(0,{#1-.5})}, yscale=-1]
% 1 pour >, -1 pour <, 0 pour rien
\foreach \it in #2 {
	\ifthenelse{ \value{col}=#1 }{
		\stepcounter{row};
		\setcounter{col}{1};
	}{}
	\ifthenelse{ \it=0 }{
	}{
		\draw[line width=2pt] ({\value{col}-\it*.15},\value{row}+.2) --++ ({\it*.3},-.2) --++ ({-\it*.3},-.2);
	}
	\stepcounter{col};
}

\end{scope}
}

\newcommand\inequalV[2]{
% Placement des symboles verticaux
\setcounter{col}{0}
\setcounter{row}{0}
\begin{scope}[shift={(.5,{#1-1})}, yscale=-1]
% 1 pour >, -1 pour <, 0 pour rien
\foreach \it in #2 {
	\ifthenelse{ \value{col}=#1 }{
		\stepcounter{row};
		\setcounter{col}{0};
	}{}
	\ifthenelse{ \it=0 }{
	}{
		\draw[line width=2pt] ({\value{col}-.2},{\value{row}-\it*.15}) --++ (.2,{\it*.3}) --++ (.2,{-\it*.3});
	}
	\stepcounter{col};
}

\end{scope}
}



\newcommand\hexjunction[2]{
	\def\c{0.8660254*#2};	
	\IfInteger{#1}{
		\draw ({#1*60}:\c) node[rotate={90-#1*60}]{$\blacklozenge$};	
	}{}
}

\newcommand\hex[5]{

\IfSubStr{#4}{?}{
    \StrBefore{#4}{?}[\nodeval];
    \StrBehind{#4}{?}[\junctions];
}{
	\def\nodeval{#4};
	\def\junctions{};
}

\IfEndWith{\nodeval}{S}
{
    \StrBefore{\nodeval}{S}[\nodevalb];
    \def\issol{1};
}
{
	\def\nodevalb{\nodeval};
	\def\issol{0}	;
}


\IfEndWith{\nodevalb}{I}
{
    \StrBefore{\nodevalb}{I}[\tag];
    \def\isstart{1};
}
{
	\def\tag{\nodevalb};
	\def\isstart{0}	;
}

\begin{scope}[shift={(#1,#2)}]
\ifthenelse{\equal{\tag}{B}}{
    \fill[black](30:#3) -- (90:#3) -- (150:#3) -- (210:#3) -- (270:#3) -- (330:#3) -- cycle;
}{
    \ifthenelse{\issol = 1}{
        \draw[fill=black!10] (30:#3) -- (90:#3) -- (150:#3) -- (210:#3) -- (270:#3) -- (330:#3) -- cycle;
    }{
        \draw (30:#3) -- (90:#3) -- (150:#3) -- (210:#3) -- (270:#3) -- (330:#3) -- cycle;
    }
    \ifthenelse{ \isstart = 1 }{
	    \def\w{{#3*.9}}
        \draw (30:\w) -- (90:\w) -- (150:\w) -- (210:\w) -- (270:\w) -- (330:\w) -- cycle;
    }{}
	\ifthenelse{#5 = 1}{% affichage cor
	    \draw (0,0) node{\tag};
	}{
	    \ifthenelse{\issol = 1}{
            \draw (0,0) node{\tag};
	    }{}
	}
}

\StrLen{\junctions}[\lenjunction]
\ifthenelse{\lenjunction = 1}{
    \hexjunction{\junctions}{#3};

}{
	\ifthenelse{\lenjunction = 2}{
	    \StrSplit{\junctions}{1}{\first}{\second}
	    \hexjunction{\first}{#3};
	    \hexjunction{\second}{#3};
	}{}
}
\end{scope}
}


\newcommand\hexGrid[3]{
\def\c{0.8660254};
\def\w{#1};
\setcounter{size}{#1}
\setcounter{col}{0}
\setcounter{row}{0}
\begin{scope}[yscale=-1]
\foreach \it in #2 {
	\ifthenelse{\value{row}<\w}{
		\hex{ {(\value{col}*2-\value{row})*\c} }{ \value{row}*1.5 }{1}{\it}{#3}
	}{
		\hex{ {(\value{col}*2+ \value{row}-\w-(\w-2))*\c} }{ \value{row}*1.5 }{1}{\it}{#3}
	}
	\stepcounter{col};	
	\ifthenelse{ \value{col} = \value{size} }{
		\stepcounter{row};
		\ifthenelse{\value{row}<\w}{
			\stepcounter{size};
		}{
		    \addtocounter{size}{-1};
		}
		
		\setcounter{col}{0};
	}{}
}
\end{scope}
}

\newcommand{\nurikabe}[4]{
    \plainGrid{#1}
    \draw[line width=1pt] (0,0) rectangle (#1,#1);
	\setcounter{col}{0}
	\setcounter{row}{0}
	\begin{scope}[shift={(0,#1)}, yscale=-1]
	\begin{scope}[shift={(.5,.5)}]
	\foreach \it in #2 {
		\ifthenelse{ \value{col}= #1 }{
			\stepcounter{row};
			\setcounter{col}{0};
		}{}
		\ifthenelse{\equal{#4}{1}}{
		    \ifthenelse{ \equal{\it}{b} }{
		        \fill[black!80] ({\value{col}-.5},{\value{row}-.5}) rectangle ++ (1,1);
		    }{}
		}{}
		\ifthenelse{ \equal{\it}{b} }{}
		{
		    \node[anchor=center, scale=#3] at (\value{col},\value{row}) {\it};
		}
		\stepcounter{col};
	}
	\end{scope}
	\end{scope}
}
